% 第5章图：无弦四环没有单纯顶点；补弦{1,3}后按(2,1,3,4)无填充删除。
\begin{tikzpicture}[x=1mm,y=1mm,>=Stealth,
  v/.style={circle,draw=FigureStroke,fill=FigureBlueFill,line width=.7pt,minimum size=5.4mm,inner sep=0pt},
  cur/.style={v,fill=FigureYellowFill,draw=FigureYellow},
  e/.style={FigureStroke,line width=.7pt},
  miss/.style={FigureRed,line width=.8pt,dashed},
  sc/.style={font=\figurefont\footnotesize,text=FigureInk},
  note/.style={font=\FigureNoteFont\fontsize{7.5}{9}\selectfont,text=FigureMuted}]
\begin{scope}
 \node[v] (a1) at (0,14) {$1$}; \node[v] (a2) at (14,14) {$2$}; \node[v] (a3) at (14,0) {$3$}; \node[v] (a4) at (0,0) {$4$};
 \draw[e] (a1)--(a2)--(a3)--(a4)--(a1);
 \draw[miss] (a1)--(a3); \draw[miss] (a2)--(a4);
 \node[sc] at (7,-10) {(a) 无弦四环};
 \node[note] at (7,-15) {没有单纯顶点};
\end{scope}
\draw[FigureGrid,line width=.6pt] (24,-16)--(24,20);
\begin{scope}[xshift=32mm]
 \node[v] (b1) at (0,14) {$1$}; \node[cur] (b2) at (14,14) {$2$}; \node[v] (b3) at (14,0) {$3$}; \node[v] (b4) at (0,0) {$4$};
 \draw[e] (b1)--(b2)--(b3)--(b4)--(b1);
 \draw[FigureBlue,line width=1.1pt] (b1)--(b3);
 \node[note] at (7,-6) {删$2$};
\end{scope}
\begin{scope}[xshift=55.5mm]
 \node[cur] (c1) at (0,14) {$1$}; \node[v] (c3) at (14,0) {$3$}; \node[v] (c4) at (0,0) {$4$};
 \draw[e] (c3)--(c4)--(c1);
 \draw[FigureBlue,line width=1.1pt] (c1)--(c3);
 \node[note] at (7,-6) {删$1$};
\end{scope}
\begin{scope}[xshift=77mm]
 \node[cur] (d3) at (11,0) {$3$}; \node[v] (d4) at (0,0) {$4$};
 \draw[e] (d3)--(d4);
 \node[note] at (5.5,-6) {删$3$};
\end{scope}
\begin{scope}[xshift=98.5mm]
 \node[cur] (e4) at (0,0) {$4$};
 \node[note] at (0,-6) {删$4$};
\end{scope}
\node[sc] at (66,-10) {(b) 补弦$\{1,3\}$后按$(2,1,3,4)$删除};
\node[note] at (66,-15) {每步都无需补边};
\end{tikzpicture}
